% !TEX program = xelatex
\documentclass[11pt,a4paper]{article}
\usepackage[UTF8]{ctex}
\usepackage{amsmath,amssymb,amsfonts}
\usepackage{booktabs}
\usepackage{graphicx}
\usepackage{hyperref}
\usepackage{geometry}
\usepackage{multirow}
\usepackage{caption}
\usepackage{subcaption}
\geometry{margin=2.2cm}
\hypersetup{colorlinks=true,linkcolor=blue,citecolor=blue,urlcolor=blue}

\title{\textbf{Adversarial Semantic Entropy (ASE) 实验报告}\\[0.4em]
\large 基于答案级语义碎裂的数学推理不确定性检测\\
\small Minerva / MATH-500 / GSM8K 三数据集对比}
\author{实验汇报稿}
\date{\today}

\begin{document}
\maketitle

\begin{abstract}
本报告在三个数学推理子集（Minerva、MATH-500、GSM8K，各约 200 题）上评估
\textbf{Adversarial Semantic Entropy (ASE)} 用于检测模型最终答案是否错误的能力。
核心思想是：在文本改写（text adversarial）与权重扰动（weight adversarial）下
\textbf{自由生成}多个语义答案，对答案做等价聚类后计算\textbf{答案簇分布的不确定性}，
而非对固定答案 $a_0$ 做 teacher-forcing 似然打分。
实验表明，联合文本--权重碎裂指标 \textbf{TW-ASE} 显著优于 TokUR 权重 epistemic baseline
（Clean AUROC：TW-ASE 0.808/0.875/0.967 \emph{vs.} TokUR EU 0.574/0.580/0.520；
另报告 AUPRC 与 TokUR 式 ACC*，见表~\ref{tab:main}）；
结构化 token/process 特征可作为小幅辅助，但难以替代答案级 ASE。
\end{abstract}

\section{研究动机}

大模型数学推理中，token 级概率往往混淆\textbf{表达不确定性}与\textbf{语义不确定性}。
Semantic Entropy 等工作指出：应将多次生成的\textbf{最终答案}归并为语义簇，再计算簇分布熵。
本实验进一步引入两类\textbf{对抗式扰动}以诱导答案分布：
\begin{itemize}
  \item \textbf{Text-side}：对题面做语义等价的 API 改写（$N{=}8$ 个 rephrase），每次自由生成完整解答；
  \item \textbf{Weight-side}：对 $q/k$ 投影施加低秩高斯权重扰动（$M{=}8$ 个 seed），在原题上重新生成。
\end{itemize}
标签 $y$ 定义为模型\textbf{基线生成} $a_0$ 是否与参考答案等价错误；不确定性分数 $s$ 越高，表示越可能错误。

\section{方法}

\subsection{单题实验流程}

对每道题 $q$，参考答案 $a^\ast$：
\begin{enumerate}
  \item \textbf{Clean 生成}：得到基线答案 $a_0$（仅用于打标签，\emph{不作为} UQ 分数）；
  \item \textbf{Text 分支}：$N{=}8$ 个 rephrase 题面 $q^{(i)}_{\mathrm{text}}$，各自自由生成，抽取最终答案；
  \item \textbf{Weight 分支}：$M{=}8$ 组权重扰动 $\theta^{(j)}$，在原题 $q$ 上生成；
  \item \textbf{联合分支}：合并 $N{+}M$ 个答案计算 TW-ASE。
\end{enumerate}
每题共 $1{+}8{+}8{=}17$ 次生成（greedy 解码，$\texttt{temperature}{=}0$）。

\subsection{数学答案聚类}

对答案集合 $\{a_k\}_{k=1}^{n}$，先用规则抽取 boxed 最终答案并归一化，再定义等价关系
$a_i \sim a_j \Leftrightarrow \texttt{math\_equal}(\tilde a_i, \tilde a_j)$。
用并查集得到 $K$ 个语义簇 $\{C_1,\ldots,C_K\}$，簇质量
\begin{equation}
  p_k = \frac{|C_k|}{n}, \quad k=1,\ldots,K.
\end{equation}

\subsection{Adversarial Semantic Entropy (ASE)}

对答案 multiset 计算（实现中主分数取 $U$）：
\begin{align}
  p_{\max} &= \max_k p_k, \\
  U_{\mathrm{ASE}} &= 1 - p_{\max} \label{eq:U} \\
  H &= -\sum_{k=1}^{K} p_k \log p_k, \\
  H_{\mathrm{norm}} &= \begin{cases}
    H / \log K, & K > 1, \\
    0, & K = 1.
  \end{cases}
\end{align}
\textbf{解释}：$U_{\mathrm{ASE}}$ 衡量``最大共识簇''之外的质量占比；
$U$ 越大，说明扰动下答案越分裂（fragmentation 越强），越可疑。
当所有生成落入同一语义簇时，$U{=}0$（稳定，无论对错）。

\paragraph{三分支定义.}
设 $\mathcal{A}_{\mathrm{T}}$、$\mathcal{A}_{\mathrm{W}}$ 分别为 text / weight 分支答案集合：
\begin{align}
  \mathrm{T\text{-}ASE} &= U(\mathcal{A}_{\mathrm{T}}), \\
  \mathrm{W\text{-}ASE} &= U(\mathcal{A}_{\mathrm{W}}), \\
  \mathrm{TW\text{-}ASE} &= U(\mathcal{A}_{\mathrm{T}} \cup \mathcal{A}_{\mathrm{W}}).
\end{align}
簇数量 $K$ 记为 \texttt{T\_num\_clusters}、\texttt{W\_num\_clusters}、\texttt{TW\_num\_clusters}，亦可单独作为检测分数。

\subsection{权重扰动（Weight Adversarial）}

对 Transformer 的 $q/k$ 投影施加\textbf{低秩高斯扰动}：
$\Delta W = \sigma \cdot (U V^\top)$，$\mathrm{rank}{=}4$，$\sigma{=}0.03$，
8 个随机种子独立采样。扰动后\textbf{整题重新生成}，记录完整 token trace（非 teacher-forcing）。

\subsection{TokUR Baseline（对比方法）}

作为\textbf{基线}，我们实现 TokUR 风格的\textbf{权重 epistemic uncertainty}：
固定基线答案 $a_0$，对 prompt + $a_0$ 做 teacher-forcing（\emph{非}自由生成），
在 $q/k$ 投影上施加 $K{=}8$ 次低秩高斯权重扰动（$\sigma{=}0.1$，$\mathrm{rank}{=}8$）。
对每个答案 token 位置 $t$，记 $K$ 次扰动下的 log-probability 为 $\{\ell_k(t)\}_{k=1}^K$，定义
\begin{equation}
  \mathrm{EU}^{W}(t) = \mathrm{Var}_{k=1..K}\big[\ell_k(t)\big],
\end{equation}
TokUR 分数取 token 级 EU 之和（与实现中 \texttt{tokur\_eu\_sum} 一致）：
\begin{equation}
  S_{\mathrm{TokUR}} = \sum_{t \in \text{answer span}} \mathrm{EU}^{W}(t).
\end{equation}
\textbf{与 ASE 的本质区别}：TokUR 问``固定 $a_0$ 在权重扰动下有多不确定'';
ASE 问``扰动后\textbf{自由生成}的语义答案是否分裂''。前者是 token-level epistemic proxy，
后者是 answer-level semantic fragmentation。

\subsection{Token / Process 辅助特征（简述）}

在 weight 分支上额外提取结构化特征，例如：
\begin{itemize}
  \item \texttt{W\_math\_token\_flip\_max}：数学 token 翻转不稳定度；
  \item \texttt{W\_alternative\_answer\_mass\_topk}：top-$k$ 中竞争答案簇质量；
  \item \texttt{final\_answer\_equiv\_last\_equation}：最终答案与末行公式是否一致（越低越可疑）。
\end{itemize}
主表仅报告与 TW-ASE 的\textbf{两特征 logistic 融合}（表~\ref{tab:fusion}），不作大规模特征工程。

\subsection{Clean Label 协议}

原始 grader 使用 \texttt{math\_equal}；\textbf{Clean label} 使用增强的
\texttt{math\_equal\_clean}（分数/矩阵/LaTeX 归一化 + 相对数值容差）。
若仅 clean 判对而原始判错，记为 \textbf{relabeled}。
Minerva 另有人工核验的 14 题 \texttt{MANUAL\_RELABEL}。
检测指标（TokUR 协议，$y{=}0$ 表示正确、分数越高越可疑）：
\textbf{AUROC}、\textbf{AUPRC}，以及 \textbf{ACC*}（按不确定性分数降序取 Top-50\% 样本的正确率）。
主结果见表~\ref{tab:main}；完整 ASE 变体见表~\ref{tab:extended}。

\section{实验设置}

\begin{table}[h]
  \centering
  \caption{实验配置摘要}
  \label{tab:setup}
  \begin{tabular}{ll}
    \toprule
    项目 & 配置 \\
    \midrule
    模型 & TFB-Qwen2.5-3B-Instruct \\
    数据集 & Minerva / MATH-500 / GSM8K（各 $\leq 200$ 题） \\
    Text rephrase & $N{=}8$（API 缓存于 \texttt{qaac\_api\_bench}) \\
    Weight perturb (ASE) & $M{=}8$，$\sigma{=}0.03$，rank$=4$，seeds 42--49 \\
    TokUR baseline & teacher-force $a_0$，$K{=}8$，$\sigma{=}0.1$，rank$=8$ \\
    解码 & greedy，\texttt{max\_new\_tokens}$=2048$ \\
    评估 & Clean-label AUROC / AUPRC / ACC*（越高越好） \\
    \bottomrule
  \end{tabular}
\end{table}

\section{实验结果}

\subsection{数据集与标签统计}

\begin{table}[h]
  \centering
  \caption{三数据集样本量与 Clean 标签统计}
  \label{tab:label}
  \begin{tabular}{lrrrr}
    \toprule
    数据集 & $N$ & Clean 准确率 & 错误数 & Relabeled \\
    \midrule
    Minerva   & 197 & 36.5\% & 125 & 22 \\
    MATH-500  & 200 & 65.5\% & 69  & 4  \\
    GSM8K     & 200 & 85.5\% & 29  & 0  \\
    \bottomrule
  \end{tabular}
\end{table}

Minerva 难度最高、标签噪声最大（22 题经 canonical grader 修正）；
GSM8K 最简单，clean 与 original 标签完全一致（relabeled$=0$）。

\subsection{表 1：Answer-level Fragmentation 主结果}

\begin{table}[t]
\centering
\caption{Clean-label detection on three math benchmarks (AUROC / AUPRC / ACC*). ACC* denotes Top-50\% accuracy.}
\label{tab:main}
\small
\begin{tabular}{l|ccc|ccc|ccc|ccc}
\toprule
 & \multicolumn{3}{c}{Minerva} & \multicolumn{3}{c}{MATH-500} & \multicolumn{3}{c}{GSM8K} & \multicolumn{3}{c}{DeepScaler} \\
\cmidrule(lr){2-4} \cmidrule(lr){5-7} \cmidrule(lr){8-10} \cmidrule(lr){11-13}
Method & AUROC & AUPRC & ACC* & AUROC & AUPRC & ACC* & AUROC & AUPRC & ACC* & AUROC & AUPRC & ACC* \\
\midrule
TokUR EU & 0.758 & 0.875 & 0.500 & 0.801 & 0.659 & 0.850 & 0.786 & 0.474 & 0.951 & 0.712 & 0.782 & 0.535 \\ \midrule
PANDA & 0.822 & 0.900 & 0.575 & 0.904 & 0.856 & 0.935 & 0.961 & 0.839 & 0.995 & 0.810 & 0.858 & 0.595 \\
F\_resp & 0.815 & 0.887 & 0.567 & 0.896 & 0.834 & 0.935 & 0.955 & 0.807 & 0.995 & 0.806 & 0.845 & 0.595 \\
D\_ans & 0.772 & 0.847 & 0.530 & 0.808 & 0.661 & 0.865 & 0.778 & 0.435 & 0.960 & 0.765 & 0.803 & 0.565 \\
\bottomrule
\end{tabular}
\end{table}


\textbf{主要观察：}
\begin{itemize}
  \item \textbf{TW-ASE} 在 AUROC/AUPRC 上全面优于 TokUR EU（AUROC 平均提升约 $+0.27$），验证 answer-level fragmentation 相对 teacher-forcing epistemic 的优势；
  \item TokUR EU 的 AUROC 仅 $\approx 0.52$--$0.58$，接近随机；但其 ACC* 在 MATH-500/GSM8K 上较高（0.66/0.85），反映 Top-50\% 选择性拒答与 AUROC 排序能力并不等价；
  \item ASE 的 ACC* 普遍低于 TokUR（如 Minerva TW-ASE ACC*$=0.11$），因高 $U$ 样本中错例占比高，Top-50\% 未必对应高准确率子集；
  \item MATH-500 上 T-ASE AUROC $>$ W-ASE；Minerva 上 W-ASE 略强于 T-ASE；
  \item GSM8K 上 TW-ASE AUROC 极高（0.967），与任务较易、错例常呈答案分裂有关（见讨论）。
\end{itemize}

\subsection{扩展表：含簇数量变体}

\input{../tables/table_extended.tex}

\subsection{表 2：Token / Process 辅助信号}

\begin{table}[h]
  \centering
  \caption{Token/process 辅助特征的 Clean AUROC}
  \label{tab:token}
  \begin{tabular}{lccc}
    \toprule
    方法 & Minerva & MATH-500 & GSM8K \\
    \midrule
    W\_math\_token\_flip\_max       & 0.748 & 0.753 & 0.678 \\
    W\_alt.\ answer\_mass (top-$k$) & 0.746 & 0.802 & 0.717 \\
    W\_confident\_fragmentation     & 0.745 & 0.814 & 0.717 \\
    W\_formula\_skeleton\_entropy   & 0.693 & 0.726 & 0.595 \\
    final\_ans.\ $\equiv$ last eq.  & 0.590 & 0.532 & 0.580 \\
    \midrule
    \textbf{TW-ASE}                 & \textbf{0.808} & \textbf{0.873} & \textbf{0.967} \\
    \bottomrule
  \end{tabular}
\end{table}

所有 token/process 指标均\textbf{低于}同数据集 TW-ASE，支持``答案级碎裂为主信号、token 为辅助''的叙事。

\subsection{表 3：Minimal Add-on 融合}

\begin{table}[h]
  \centering
  \caption{以 TW-ASE 为主、叠加 1--2 个辅助特征的 Logistic 融合 AUROC}
  \label{tab:fusion}
  \begin{tabular}{lccc}
    \toprule
    模型 & Minerva & MATH-500 & GSM8K \\
    \midrule
    TW-ASE                              & 0.808 & 0.873 & 0.967 \\
    TW-ASE + alt.\ answer mass          & 0.816 & 0.874 & 0.967 \\
    TW-ASE + final $\equiv$ last eq.    & 0.820 & 0.874 & 0.968 \\
    TW-ASE + both                       & 0.831 & 0.876 & 0.967 \\
    \bottomrule
  \end{tabular}
\end{table}

Minerva 上融合增益最明显（$+0.02\sim0.03$）；
MATH-500 / GSM8K 上 TW-ASE 已接近饱和，融合收益有限。

\section{讨论}

\paragraph{1. 为何放弃 teacher-forcing / TokUR 主线？}
TokUR EU 固定 $a_0$ 做 teacher-forcing，衡量权重扰动下 token log-prob 的方差；
本实验测得 AUROC 仅 0.52--0.58（表~\ref{tab:main}），接近随机。
ASE 则在扰动下\textbf{重新生成}并观察语义答案簇是否分裂，TW-ASE 达 0.81--0.97。
这说明数学 UQ 的关键信号在\textbf{语义答案分布}，而非固定轨迹上的 token epistemic 波动。

\paragraph{2. Minerva 标签清洗的重要性.}
22 题 relabel 将准确率从 25.4\% 提升至 36.5\%，避免格式等价误判系统性压低 AUROC。
汇报时应同时说明 \textbf{Original vs Clean} 标签差异。

\paragraph{3. GSM8K 上 TW-ASE $=$ 0.967 是否``夸张''？}
清洗流程已执行，但 \textbf{relabeled$=0$}，Clean 与 Original AUROC 相同。
该子集 Clean 准确率 85.5\%，29 个错例中 24 个 TW-ASE $\geq 0.3$，仅 5 个 stable-wrong；
Bootstrap 95\% CI 约为 $[0.94,\,0.99]$。
因此高 AUROC 主要来自\textbf{任务较易 + 错例多呈答案碎裂}，而非标签被洗高。
同时 W-ASE 仅 0.723，说明并非所有指标都异常偏高。

\paragraph{4. 与 Semantic Entropy 的关系.}
本工作可视为在\textit{对抗扰动诱导的答案分布}上计算语义熵/碎裂度：
$H$ 与 $H_{\mathrm{norm}}$ 已记录，但主表统一采用 $U_{\mathrm{ASE}}=1-p_{\max}$，因其对``单簇主导''更直观。

\section{结论}

\begin{enumerate}
  \item \textbf{Answer-level fragmentation (TW-ASE)} 显著优于 TokUR EU baseline，是跨三数据集的最强检测信号；
  \item \textbf{Token/process} 特征有判别力，但应作为轻量辅助而非主线；
  \item \textbf{Clean label} 在 Minerva 等复杂题型上不可或缺；GSM8K 需结合难度与 CI 解读高 AUROC；
  \item 后续可扩展 DeepScaler 等更难数据集，并报告 Original/Clean 双轨指标与 stable-wrong 案例分析。
\end{enumerate}

\appendix
\section{核心公式速查}

\begin{align*}
  &\text{簇质量: } p_k = |C_k|/n, \quad
  U_{\mathrm{ASE}} = 1 - \max_k p_k, \quad
  H = -\textstyle\sum_k p_k \log p_k \\
  &\mathrm{T\text{-}ASE} = U(\mathcal{A}_{\mathrm{T}}),\;
   \mathrm{W\text{-}ASE} = U(\mathcal{A}_{\mathrm{W}}),\;
   \mathrm{TW\text{-}ASE} = U(\mathcal{A}_{\mathrm{T}} \cup \mathcal{A}_{\mathrm{W}})
\end{align*}

\section{复现路径}

\begin{itemize}
  \item 原始数据：\texttt{/HDDDATA/phx/PRS/outputs/ase\_full/\{dataset\}/raw\_runs/}
  \item 论文表格（AUROC/AUPRC/ACC*）：\texttt{python -m prs.ase.analyze\_paper\_tables --out-dir ... --paper-dir paper/}
  \item 单数据集 AUROC 表：\texttt{python -m prs.ase.analyze\_main\_tables --dataset <name>}
  \item TokUR baseline：\texttt{python -m prs.ase.score\_tokur\_baseline --dataset <name>}
  \item 生成物：\texttt{paper/tables/table\_main.\{md,tex\}}, \texttt{table\_extended.\{md,tex\}}, \texttt{results.json}
\end{itemize}

\end{document}
